% SEGMENT OLP-0088-S01
% Part: first-order-logic
% Chapter: natural-deduction
% Section: derivations

\documentclass[../../../include/open-logic-section]{subfiles}

\begin{document}

\iftag{FOL}
      {\olfileid{fol}{ntd}{der}}
      {\olfileid{pl}{ntd}{der}}

\olsection{\usetoken{P}{derivation}}

\begin{explain}
எடுகோள் என்றால் என்ன என்று கூறி, அனுமான விதிகளையும் கொடுத்துள்ளோம். இயல்பான
வருவித்தலில் உள்ள !!^{derivation}s இவற்றிலிருந்து தொகுத்தறிதல் முறையில்
உருவாக்கப்படுகின்றன: ஒவ்வொரு !!{derivation} உம் தனியே ஓர் எடுகோளாகவோ,
சரியான ஓர் அனுமானத்தைத் தொடர்ந்து வரும் ஒன்று, இரண்டு அல்லது மூன்று
!!{derivation}s ஆகவோ இருக்கும்.
\end{explain}

% SEGMENT OLP-0088-S02
\begin{defn}[!!^{derivation}]
எடுகோள்கள்~$\Gamma$ இலிருந்து !!a{sentence}~$!A$ இன்
\Article{derivation} \emph{!!{derivation}} என்பது பின்வரும் நிபந்தனைகளை
நிறைவுறுத்தும் !!{sentence}s[gen] ஒரு முடிவுறு மரவுரு:
\begin{enumerate}
\item மரவுருவின் மிக மேலுள்ள !!{sentence}s, $\Gamma$ வில் உள்ளன அல்லது
  மரவுருவிலுள்ள ஓர் அனுமானத்தால் !!{discharged} ஆகின்றன.
\item மரவுருவின் மிகக் கீழுள்ள !!{sentence} ஆனது~$!A$.
\item மரவுருவின் கீழுள்ள !!{sentence}~$!A$ ஐத் தவிர அதிலுள்ள ஒவ்வொரு
  !!{sentence} உம், அதற்கு நேரடியாகக் கீழே முடிவு அமைந்திருக்கும் அனுமான
  விதியின் சரியான பயன்பாடு ஒன்றின் முன்கோளாகும்.
\end{enumerate}
அப்போது $!A$ ஐ அந்த !!{derivation}[gen] \emph{முடிவு} என்றும், $\Gamma$ வை
அதன் !!{undischarged} எடுகோள்கள் என்றும் சொல்கிறோம்.

$\Gamma$ இலிருந்து $!A$ இன் !!a{derivation} இருந்தால், $!A$ ஆனது
$\Gamma$ இலிருந்து \emph{!!{derivable}} என்கிறோம்; குறியீட்டில்:
$\Gamma \Proves !A$. ஒவ்வோர் எடுகோளும் !!{discharged} ஆகும் $!A$ இன்
!!a{derivation} இருந்தால், $\Proves !A$ என்று எழுதுகிறோம்.
\end{defn}

% SEGMENT OLP-0088-S03
\begin{ex}
ஒவ்வோர் எடுகோளும் தனியே !!a{derivation}. ஆகவே, எடுத்துக்காட்டாக, $!A$ தனியே
!!a{derivation}; $!B$ உம் தனியே அப்படியே. இவற்றிற்கு, எடுத்துக்காட்டாக,
$\Intro{\land}$ விதியைப் பயன்படுத்தி ஒரு புதிய !!{derivation} ஐப் பெறலாம்:
\begin{prooftree}
\AxiomC{$!A$}
\AxiomC{$!B$}
\RightLabel{\Intro{\land}}
\BinaryInfC{$!A \land !B$}
\end{prooftree}
இந்த விதிகள் பொதுவானவையாக இருக்க நோக்கமுடையவை: அதிலுள்ள $!A$, $!B$ ஐ எந்த
!!{sentence}s ஆலும், எடுத்துக்காட்டாக $!C$, $!D$ ஆலும் மாற்றலாம். அப்போது
முடிவு $!C \land !D$ ஆகும்; எனவே
\begin{prooftree}
\AxiomC{$!C$}
\AxiomC{$!D$}
\RightLabel{\Intro{\land}}
\BinaryInfC{$!C \land !D$}
\end{prooftree}
ஒரு சரியான !!{derivation}. நிச்சயமாக, எடுகோள்களை இடமாற்றவும் முடியும்; அப்போது
$!D$ ஆனது~$!A$ இன் பங்கையும் $!C$ ஆனது~$!B$ இன் பங்கையும் வகிக்கும். ஆகவே,
\begin{prooftree}
\AxiomC{$!D$}
\AxiomC{$!C$}
\RightLabel{\Intro{\land}}
\BinaryInfC{$!D \land !C$}
\end{prooftree}
உம் ஒரு சரியான !!{derivation}.

% SEGMENT OLP-0088-S04
இப்போது மற்றொரு விதியை, எடுத்துக்காட்டாக $\Intro{\lif}$ ஐப் பயன்படுத்தலாம்.
இது ஒரு நிபந்தனைவாக்கியத்தை முடிவாகக் கொள்ள அனுமதிப்பதுடன், அந்த
நிபந்தனைவாக்கியத்தின் முன்னிலைக்கு ஒன்றான எந்த எடுகோளையும் !!{discharge}
செய்ய அனுமதிக்கிறது. ஆகவே பின்வரும் இரண்டும் சரியான !!{derivation}s ஆகும்:
\begin{prooftree}
\AxiomC{$\Discharge{!C}{1}$}
\AxiomC{$!D$}
\RightLabel{\Intro{\land}}
\BinaryInfC{$!C \land !D$}
\DischargeRule{\Intro{\lif}}{1}
\UnaryInfC{$!C \lif (!C \land !D)$}
\DisplayProof\bottomAlignProof
\AxiomC{$!C$}
\AxiomC{$\Discharge{!D}{1}$}
\RightLabel{\Intro{\land}}
\BinaryInfC{$!C \land !D$}
\DischargeRule{\Intro{\lif}}{1}
\UnaryInfC{$!D \lif (!C \land !D)$}
\end{prooftree}
அவை முறையே, $!D \Proves !C \lif (!C \land !D)$ என்றும்
$!C \Proves !D \lif (!C \land !D)$ என்றும் காட்டுகின்றன.

% SEGMENT OLP-0088-S05
எடுகோள்களை விடுவித்தல் ஓர் அனுமதியே அன்றி, கட்டாயம் அல்ல என்பதை நினைவுகூர்க:
எடுகோள்களை விடுவிக்க வேண்டியதில்லை. குறிப்பாக, விடுவிக்க வேண்டிய எடுகோள்கள்
!!{derivation}[loc] இல்லாதபோதும் ஒரு விதியைப் பயன்படுத்தலாம். எடுத்துக்காட்டாக,
!!{discharged} செய்யப்பட வேண்டிய எடுகோள்~$!A$ எதுவும் இல்லாதபோதும்,
பின்வருவது முறையானது:
\begin{prooftree}
  \AxiomC{$!B$}
  \DischargeRule{\Intro{\lif}}{1}
  \UnaryInfC{$!A \lif !B$}
\end{prooftree}
\end{ex}

\end{document}
